% ===================================================================
%  সারণি নং 2 — মানবদেহ। — অবকাশ। — খেলা।
%
%  Ceci n'est PAS une transcription : c'est une traduction, faite en
%  2026, en bengali d'aujourd'hui. Voir texto/bn/10-sarani-01.tex
%  pour la consigne, qui vaut pour les seize fichiers.
% ===================================================================

%%K t02-apar-1 apar
\VUsaut{4.0mm}
\VUtitre{15.0pt}{60}{সারণি নং 2}
\VUsaut{2.0mm}
\VUtitre{11.4pt}{0}{মানবদেহ। --- অবকাশ।}
\VUtitre{11.4pt}{0}{খেলা।}

%%K t02-tit-0 sub
\VUtitre{13.2pt}{0}{মানবদেহ।}
\VUsaut{2.4mm}
\VUcentre{10.2pt}{0}{\textit{(প্রাকৃতিক বিজ্ঞানের একটি সহজ পাঠ।)}}
\VUsaut{3.0mm}

%%K t02-01-1 p
1. --- মানবদেহের প্রধান অংশগুলি হল: \VUgras{মাথা}, \VUgras{ধড়} ও
\VUgras{অঙ্গপ্রত্যঙ্গ}।

%%K t02-tit-1 sub
\VUpk{10.2pt}{40}{মাথা।}
\VUsaut{3.0mm}

%%K t02-02-1 p
2. --- মাথার দুটি অংশ: \VUgras{করোটি}, যার ভিতরে \VUgras{মস্তিষ্ক} থাকে ও
যাকে \VUgras{চুল} ঢেকে রাখে, আর \VUgras{মুখমণ্ডল}।

%%K t02-03-1 p
3. --- আমাদের ছবিতে করোটিও দেখা যায় না, মস্তিষ্কও নয়; কিন্তু মুখমণ্ডলের
গুরুত্বপূর্ণ অংশগুলি খুব স্পষ্ট দেখা যায়।

%%K t02-04-1 p
4. --- এই যে প্রথমে \VUgras{ললাট}, যা প্রবল বুদ্ধি ও সদা জাগ্রত মনের পরিচয়
দেয়। \VUgras{রগ} দুটি খুব খোলা। \VUgras{চোখ} চঞ্চল ও প্রশস্ত। একটি ঘন
\VUgras{ভুরু} ও লম্বা \VUgras{পলক} তাকে রক্ষা করে।

%%K t02-05-1 p
5. --- এই যে \VUgras{নাক} ও \VUgras{নাসারন্ধ্র}। নাক দিয়ে আমরা গন্ধ নিই ও
শ্বাস নিই। নাকের দুই পাশে \VUgras{গাল}, আর তার পরে \VUgras{কান}। মুখমণ্ডলের
নিচের অংশ \VUgras{মুখ} ও \VUgras{চিবুক} দিয়ে গড়া।

%%K t02-06-1 p
6. --- সেগুলি ভালো করে দেখা যায় না, কারণ \VUgras{গোঁফ} ও \VUgras{দাড়ি}
সেগুলিকে ঢেকে রাখে। তবু আমরা জানি যে মুখে দুটি \VUgras{ঠোঁট},
\VUgras{উপরের চোয়াল} ও \VUgras{নিচের চোয়াল} তার \VUgras{দাঁতসহ}, এবং
\VUgras{জিভ} থাকে। মুখ দিয়ে আমরা কথা বলি ও খাই। জিভ দিয়ে আমরা খাবারের
স্বাদ নিই।

%%K t02-07-1 p
7. --- চোখ, কান, নাক ও মুখ, আঙুলের সঙ্গে, পাঁচ ইন্দ্রিয়ের অঙ্গ:
\VUgras{দৃষ্টি}, \VUgras{শ্রবণ}, \VUgras{ঘ্রাণ}, \VUgras{স্বাদ} ও
\VUgras{স্পর্শ}।

%%K t02-08-1 p
8. --- এইভাবে মাথা মানবদেহের সবচেয়ে কৌতূহলজনক অংশগুলির একটি।

%%K t02-tit-2 sub
\VUsaut{4.4mm}
\VUpk{10.2pt}{40}{ধড়।}
\VUsaut{3.0mm}

%%K t02-09-1 p
9. --- \VUgras{ঘাড়} মাথাকে ধড়ের সঙ্গে জোড়ে। তার মধ্যে \VUgras{গলা} ও
\VUgras{ঘাড়ের পিছন} পড়ে।

%%K t02-10-1 p
10. --- ধড়েও অনেক অপরিহার্য অঙ্গ থাকে। \VUgras{বুকে} আছে \VUgras{ফুসফুস},
যা দিয়ে আমরা শ্বাস নিই, আর \VUgras{হৃৎপিণ্ড}, যা জীবনের কেন্দ্র।
হৃৎপিণ্ড স্পন্দন থামালে জীব মারা যায়।

%%K t02-11-1 p
11. --- \VUgras{পাঁজর} বুককে ধরে রাখে।

%%K t02-12-1 p
12. --- ধড়ের অন্য অংশগুলি হল \VUgras{পাঁজরের পাশ}, \VUgras{পিঠ},
\VUgras{কাঁধ} ও \VUgras{উদর} (পেট)। ঘুমাতে আমরা পাশ ফিরে বা চিত হয়ে শুই,
আর বোঝা কাঁধে বই।

%%K t02-13-1 p
13. --- উদরে আছে \VUgras{পাকস্থলী} ও \VUgras{অন্ত্র}। সেগুলি খাবার হজম
করতে কাজে লাগে।

%%K t02-tit-3 sub
\VUsaut{4.4mm}
\VUpk{10.2pt}{40}{অঙ্গপ্রত্যঙ্গ।}
\VUsaut{3.0mm}

%%K t02-14-1 p
14. --- আমাদের চারটি \VUgras{অঙ্গ}: দুটি \VUgras{বাহু} ও দুটি \VUgras{পা}।
বাহু দিয়ে আমরা জিনিস ধরি, ধরে রাখি, তুলি ও কুড়াই; পা দিয়ে আমরা দাঁড়াই,
হাঁটি ও দৌড়াই।

%%K t02-15-1 p
15. --- \VUgras{কাঁধ} বাহুকে দেহের সঙ্গে জোড়ে। একটি খুব শক্ত
\VUgras{পেশি}, দ্বিশির, বাহু নাড়ায়, যাকে \VUgras{কনুই}
\VUgras{বাহুমূলের নিচের অংশের} সঙ্গে জোড়ে। \VUgras{কবজি} \VUgras{হাতের}
জোড় গড়ে, যা \VUgras{আঙুল} ও \VUgras{নখে} শেষ হয়। হাতে আমরা একটি
\VUgras{শিরা} (ও একটি ধমনি) দেখতে পাই, যার মধ্যে \VUgras{রক্ত} বয়।

%%K t02-16-1 p
16. --- পায়ের অংশগুলি হল: \VUgras{ঊরু}, \VUgras{হাঁটু} ও
\VUgras{হাঁটুর পিছন}, \VUgras{ডিম}, যার \VUgras{মাংস} \VUgras{চামড়ায়}
ঢাকা, \VUgras{গোড়ালির গাঁট} ও \VUgras{পা}। পায়ের \VUgras{হাড়} খুব মোটা ও
খুব দৃঢ়। পায়ের ডগায় আছে \VUgras{পায়ের আঙুল}। অন্য অংশগুলি
\VUgras{পায়ের পাতা} ও \VUgras{গোড়ালি}।

%%K t02-17-1 p
17. --- মানবদেহের প্রায় সম্পূর্ণ বর্ণনা এই।

%%K t02-17-2 p
\textit{টীকা।} --- \textit{« আমার ... ব্যথা করছে (মাথা, ইত্যাদি) » কথাটির
সাহায্যে শিক্ষক এই সমস্ত শব্দভান্ডার ভালো বা মন্দ} \VUgras{শারীরিক
অবস্থার} \textit{দিক থেকে আবার ঝালিয়ে নিতে পারবেন।}

%%K t02-tit-4 sub
\VUsaut{5.0mm}
\VUpk{10.2pt}{40}{ব্যায়াম।}
\VUsaut{2.4mm}
\VUcentre{10.2pt}{0}{\textit{(এক ছাত্রের জবানিতে।)}}
\VUsaut{3.0mm}

%%K t02-18-1 p
18. --- নমনীয়, ক্ষিপ্র ও সুস্থ থাকতে হলে আমাদের পা ও বাহু খাটাতে হয়।
তাই আমরা খেলি এবং \VUgras{ব্যায়াম} করি।

%%K t02-19-1 p
19. --- সপ্তাহের দিনগুলিতে এই দুই অভ্যাস বিদ্যালয়ের উঠানে হয় (দেখুন
সারণি নং 1)। কিন্তু বৃহস্পতিবার ও রবিবার আমরা একটি বড় মাঠে যাই, যেখানে
দৌড়ানোর ও মজা করার জায়গা মেলে।

%%K t02-20-1 p
20. --- আমাদের শিক্ষকেরা ও আমাদের বাবা-মা কখনও কখনও সঙ্গে আসেন; কিন্তু
অভ্যাস করান \VUgras{ব্যায়াম-শিক্ষকই} \textsuperscript{(12)}।

%%K t02-21-1 p
21. --- মাঠে, প্রবেশপথের বাঁ দিকে, \VUgras{ব্যায়ামের সরঞ্জাম}
\textsuperscript{(1)} রাখা: আমরা \VUgras{মই} \textsuperscript{(2)} বেয়ে উঠি,
\VUgras{রিং} \textsuperscript{(3)} ও \VUgras{দোলনা-দণ্ডে} \textsuperscript{(4)}
উলটো ঝুলি, \VUgras{দড়ির মই} \textsuperscript{(5)} বা \VUgras{গিঁটওয়ালা
দড়ির} \textsuperscript{(6)} ডগা পর্যন্ত হাতে হাতে উঠে যাই।

%%K t02-22-1 p
22. --- এদিকে আমাদের সঙ্গীরা \VUgras{কাঠের ঘোড়ার} \textsuperscript{(7)} বা
\VUgras{সমান্তরাল দণ্ডের} \textsuperscript{(11)} উপর দিয়ে লাফায়। কেউ
\VUgras{মুষ্টিযুদ্ধ} \textsuperscript{(8)} করে বা মুখোশ ও \VUgras{ফয়েল} নিয়ে
\VUgras{অসিক্রীড়া} \textsuperscript{(9)} করে। কেউ আবার \VUgras{ডাম্বেল}
\textsuperscript{(10)} তোলে।

%%K t02-tit-5 sub
\VUsaut{4.6mm}
\VUpk{10.2pt}{40}{খেলা।}
\VUsaut{3.0mm}

%%K t02-23-1 p
23. --- ব্যায়াম করিয়েদের চারপাশে ছেলেমেয়েরা নানা রকম \VUgras{খেলা}
খেলছে। মেয়েরা তাদের \VUgras{চাকা} \textsuperscript{(14)} নিয়ে মজা করে; তারা
\VUgras{দড়ি লাফায়} \textsuperscript{(13)}, অথবা নিজেদের ভাই ও খুড়তুতো ভাইদের
\VUgras{ক্রোকে} \textsuperscript{(15)} খেলায় ডাকে। প্রতিপক্ষ যখন ভাবছে সে
অনায়াসে খিলানটা পেরিয়ে যাবে, ঠিক তখনই তারা নিজেদের \VUgras{কাঠের
মুগুর} দিয়ে তার \VUgras{বল} দূরে সরিয়ে দেয়, আর এতে তাদের ভারী আনন্দ হয়!

%%K t02-24-1 p
24. --- ছেলেরা বেশি জোরালো খেলা পছন্দ করে। তারা খুশি হয়ে \VUgras{নয় গুটি}
\textsuperscript{(16)} খেলে, যেগুলি তারা কায়দা করে \VUgras{বল} \textsuperscript{(17)}
দিয়ে ফেলে দেয়। তারা \VUgras{লাট্টু} \textsuperscript{(18)} ও
\VUgras{বিলবোকে} \textsuperscript{(19)}, এমনকি \VUgras{ব্যাঙের খেলাকেও}
\textsuperscript{(20)}, তুচ্ছ ভাবে না। কিন্তু তাদের প্রিয় খেলা \VUgras{মার্বেল}
\textsuperscript{(21)} ও \VUgras{হাতের বল} \textsuperscript{(22)}, কারণ
\VUgras{কাচঘরের} \textsuperscript{(26)} দেয়ালটি হালকা বল লোফার পক্ষে খুব
উপযোগী।

%%K t02-25-1 p
25. --- ছোট্ট জন তার \VUgras{রণপায়ে} \textsuperscript{(24)} চড়ে খুব দক্ষ এবং
নিজের ভারসাম্য বেশ রাখতে জানে। তার কাছেই কয়েকজন সঙ্গী সেই কাচঘরে ও
\VUgras{বেড়ার} পিছনে \VUgras{লুকোচুরি} \textsuperscript{(25)} খেলছে। কারও ভালো
লাগে \VUgras{চাপড়-কানামাছি} \textsuperscript{(28)} বা \VUgras{শাটল}
\textsuperscript{(29)}, যা এক \VUgras{ব্যাট} থেকে অন্যটিতে উড়ে যায়।

%%K t02-noto-1 noto
\VUnotes{143.2mm}{%
(*) শিশুদের খেলার নাম প্রায়ই একেবারে জাতীয় হয়। আমরা যেমন পেরেছি তেমনই
সেগুলিকে ইদোতে নাম দিয়েছি — কখনও যে ক্রিয়াটি তাদের চেনায় তা থেকে, কখনও
কোনো না কোনো দেশের নাম নিয়ে, যখন তা বেশি ভুল বোঝায় না। যেমন:
\VUgras{পিপের খেলা} (জার্মান ও ফরাসি) সেই \VUgras{ব্যাঙের খেলাই}
\textsuperscript{(20)} (ইংরেজি), যাতে খেলোয়াড়েরা নিজেদের গোল চাকতি ছিদ্রওয়ালা
সিন্দুকের (পিপের) উপরে হাঁ করে বসা ধাতুর ব্যাঙের মুখে ফেলার চেষ্টা করে।
\VUgras{কাঁধ-লাফ} \textsuperscript{(30)} \VUgras{ব্যাঙ-লাফ} থেকে কেবল এইটুকু
আলাদা: মাথার উপর দিয়ে লাফাতে খেলোয়াড়েরা সঙ্গীর কাঁধে হাত রাখে, অথচ
« \VUgras{ব্যাঙ-লাফে} » তারা কেবল তার পিঠে হাত রাখে। --- অথবা, কয়েকজন
খেলোয়াড় ঝুঁকে দাঁড়ায় আর অন্যেরা তাদের পিঠে হাত রেখে তাদের উপর দিয়ে
লাফায়; প্রথমদের না নড়ে ধাক্কা সইতে হয়, দ্বিতীয়দের ভারসাম্য না হারিয়ে
লাফাতে হয় (সারণিটিই দেখুন)।%
}

%%K t02-26-1 p
26. --- মাঠের মাঝখানে দুটি গাছ আছে। তার একটির নিচে সাত-আটজন ছেলে
\VUgras{কাঁধ-লাফের} \textsuperscript{(30)} (*) খেলা জমিয়েছে। এই খেলা সব দেশে
জানা নেই; কিন্তু \VUgras{ব্যাঙ-লাফ} কী, তা সবাই জানে, যা ছেলেরা এখন খেলছে
না।

%%K t02-tit-6 sub
\VUsaut{5.0mm}
\VUpk{10.2pt}{40}{খোলা হাওয়ার খেলা।}
\VUsaut{3.4mm}

%%K t02-27-1 p
27. --- \VUgras{কানামাছিতেও} \textsuperscript{(31)} খুব মজা; মেয়েরা ও ছেলেরা
পালা করে চোখে \VUgras{পট্টি} বাঁধে আর যারা ধরা দেয় সেই আনাড়িদের ধরে।

%%K t02-28-1 p
28. --- মাঠের মাঝখানে, এই যে একটি খুব ফরাসি খেলা, যদিও কিছুটা রুক্ষ: এটি
\VUgras{ভালুক} \textsuperscript{(32)}, যা \VUgras{ঘুড়ির} \textsuperscript{(33)} শান্ত
খেলার চেয়ে, এমনকি ইংরেজি খেলা \VUgras{টেনিসের} \textsuperscript{(34)} চেয়েও,
অনেক বেশি হইচইয়ের।

%%K t02-29-1 p
29. --- সেই শেষ খেলাটি বড় ছাত্রদের আনন্দ। আমাদের শিক্ষকেরা নিজেরাও তা
খুশি হয়ে খেলেন। তাতে অনেক দক্ষতা ও ক্ষিপ্রতা লাগে, আর আমার তা ভীষণ ভালো
লাগে। \VUgras{দোলনা} \textsuperscript{(35)} আমি মেয়েদের জন্য ছেড়ে দিই।

%%K t02-30-1 p
30. --- আর একটি ইংরেজি খেলাও আমার ভালো লাগে না, যা সহজেই বিপজ্জনক হয়ে
উঠতে পারে।

%%K t02-30-1b p
আমি \VUgras{ফুটবলের} \textsuperscript{(36)} কথা বলছি, যা আমার বড় ভাইয়ের আনন্দ।
আমার বেশি ভালো লাগে আমাদের পুরনো খেলা \VUgras{বন্দি-ঘাঁটি}
\textsuperscript{(38)}, ততটাই স্বাস্থ্যকর আর অনেক কম কঠোর।

%%K t02-tit-7 sub
\VUsaut{4.6mm}
\VUpk{10.2pt}{40}{সাইকেল ও বলের খেলা।}
\VUsaut{3.0mm}

%%K t02-31-1 p
31. --- মাঠের চারপাশে \VUgras{সাইকেল} \textsuperscript{(37)} চালাতেও আমার ভালো
লাগে।

%%K t02-32-1 p
32. --- আগামী বছর আমি \VUgras{ঘোড়ায় চড়া} \textsuperscript{(39)} শিখব। ঘোড়া
কী সুন্দর দেখায় যখন সে ছন্দে হাঁটে বা সুঠাম দুলকি চালে চলে, আর সরপটে বাধা
পেরিয়ে যায়!

%%K t02-33-1 p
33. --- গরমকালে আমরা \VUgras{সাঁতার} \textsuperscript{(40)} শিখি। নদীর স্বচ্ছ
ঠান্ডা জলে আমরা আনন্দে ডুব দিই ও স্নান করি। কখনও কখনও শেষে আমরা একটি
কাগজের \VUgras{বেলুন} \textsuperscript{(41)} ওড়াই, যা গরম হাওয়ায় ভরে উঠতেই
গম্ভীরভাবে উপরে উঠে যায়।

%%K t02-34-1 p
34. --- আরও অনেক খেলা আছে, কিন্তু সেগুলির নাম করতে করতে আমি কখনও ফুরাব
না, আর এই ছোট বিবরণেই দেখা যায় যে আমাদের সুন্দর মাঠে বৃহস্পতিবারগুলি
একটুও একঘেয়ে নয়।

%%K t02-34-2 p
বি. দ্র. --- \textit{শিক্ষক সহজেই এই অধ্যায়টি পূর্ণ করতে পারবেন, প্রতিটি
গুরুত্বপূর্ণ খেলার আরও বিশদ বর্ণনা দিয়ে।}
\VUsaut{6.0mm}
\VUfilet{22mm}
